% Section VII: Conclusions and future work
\section{Conclusions and Future Work}
\label{sec:conclusiones}
The poor data efficiency of a small ViT trained from scratch is not an inherent limitation of the Transformer, but the consequence of a locality bias that attention does not learn from little data. With 16 paired seeds, adding a $5\times5$ depthwise convolution over the values of each head contributes $+5.31$ points on CIFAR-10 over an equal-capacity control ($+5.64$ over the base model) with 3,744 effective parameters, improves calibration and reduces by a factor of 3.7 the GPU time needed to reach 79~\% validation accuracy, close to the final value of the base model. Over the same control, the gain recurs on CIFAR-100 ($+5.4$) and SVHN ($+2.8$); it holds with CutMix and from 119,000 to 3.6 million parameters, and it ties with LocalViT at a lower cost. A ResNet of the same size is still better, and the term closes three quarters of that gap. A global summary branch on top of the neighborhood, which completes the triad, adds $+0.61$ points over its capacity control with 48 seeds ($+0.67$ on the 32 seeds added after deciding the extension) and improves calibration, but the triad does not meet the cost criterion set in advance. Combining two similarity functions contributes at most half a point, which cannot be distinguished from capacity. The useful heterogeneity in this regime is that of inductive biases.

Four design guidelines follow from the results:
\begin{enumerate}
  \item \textbf{Avoid score fusion before the \emph{softmax}}: it multiplies the distributions; the branch with the larger scale decides and the other has no effect.
  \item \textbf{Control capacity before attributing a gain}: without a local term, the advantage of post-\emph{softmax} concatenation disappears against a wider MLP with the same parameters, and with it 48 seeds were needed to separate a 0.6-point gain from capacity.
  \item \textbf{Add a $5\times5$ local term}: initialized at zero, it can be added to an existing model without changing its initial behavior, and it costs 3.4~\% more MACs and 1.08 times the epoch time.
  \item \textbf{Combine different biases, not similarities}: the triad (neighborhood, global summary and pairwise attention) keeps them in separate terms, fused after the \emph{softmax} or added as residuals. Each term should be checked against an equal-capacity control to contribute on its own on the dataset of interest: in this study the neighborhood provides almost all of the gain ($+5.3$ over capacity), and the global summary $+0.6$ points and better calibration.
\end{enumerate}

Three lines remain as future work: measuring whether the advantage of the local term shrinks with pretraining and large datasets, where attention can learn locality by itself (up to 3.6 million parameters trained from scratch it does not shrink); integrating it into architectures for semantic segmentation and object detection; and evaluating its deployment on edge hardware, which this study does not evaluate: quantization-aware training to 8-bit integers and measurement of latency, memory and energy consumption on microcontrollers and NPUs.
